\section{Discussion}

\paragraph{Why Split-Context Provides Robustness.}
A key insight from our work is that robustness emerges from \emph{architectural isolation} rather than explicit defense mechanisms.
Full-context methods are fundamentally vulnerable because adversarial content in one evidence source can influence reasoning about unrelated sources---a fake review containing ``ANSWER: recommend'' affects judgments about metadata attributes processed in the same context.
\method's split-context design eliminates this attack vector by construction: metadata checks never see review text, and each review is processed independently.
This architectural defense requires no additional prompting and introduces no performance penalty on clean data, unlike defense prompts that make models overly conservative.

\paragraph{Theoretical Grounding.}
Our design draws on the observation that LLM accuracy decays super-linearly with context length~\cite{liu2024lost,levy2024same}.
Formally, if success probability follows $P(L) = \exp(-aL^\gamma)$ for context length $L$ with decay exponent $\gamma > 1$, then decomposing a task into $K$ independent sub-tasks of length $L/K$ yields joint success probability $P_{\text{split}} = \exp(-aL^\gamma K^{1-\gamma})$, which exceeds $P_{\text{full}}$ when $\gamma > 1$.
Empirically, $\gamma \approx 1.1$--$2$ for GPT-class models, validating the decomposition benefit.
Beyond accuracy, split-context provides a \emph{security} benefit: adversarial content in one segment cannot propagate errors to unrelated segments, as each sub-task operates in isolation.
This theoretical foundation explains why \method maintains robustness without sacrificing clean-data performance.

\paragraph{Preserving Minority Viewpoints.}
Three-valued logic serves a dual purpose in \method.
First, it handles genuine uncertainty (missing attributes, ambiguous reviews) without forcing false binary decisions.
Second, and more subtly, it preserves minority viewpoints during aggregation.
When reviews are processed independently and aggregated with explicit unknown handling, a single review noting ``the WiFi was slow'' cannot be drowned out by ten reviews that simply fail to mention WiFi.
The unknown values propagate through the Boolean logic, ensuring that absence of evidence is not treated as evidence of absence.

\paragraph{Limitations.}
\method incurs additional computational cost compared to single-pass methods: multiple LLM calls (one per evidence source plus aggregation) versus one.
This overhead is partially mitigated by caching validated scripts across items, but remains a consideration for latency-sensitive applications.
Additionally, our evaluation focuses on structured recommendation with discrete Boolean conditions; tasks requiring fuzzy matching or continuous scoring may need different aggregation mechanisms.
Finally, while our adaptation phase detects several attack patterns, sophisticated adversaries might craft attacks that evade detection---though architectural isolation still provides baseline protection.

\paragraph{Counter-Intuitive Finding.}
Adding explicit defense prompts (``watch out for adversarial content'', ``be skeptical of suspicious reviews'') to full-context baselines actually \emph{hurts} performance on clean data.
Models become overly conservative, treating legitimate reviews with strong opinions as potentially adversarial.
This suggests that prompting alone cannot solve the robustness problem---architectural changes that prevent adversarial content from reaching the reasoning context are necessary.
